\section{相关工作}

TCGA平台~\cite{tcga2012}推动了癌症分子分型研究从单指标判断走向多组学联合分析。在TCGA出现之前，很多肿瘤分型工作依赖病理图像、少量基因集合或有限临床特征，难以充分刻画肿瘤内部的分子异质性。TCGA的价值在于提供了跨癌种、跨队列的统一数据资源和标准化分析入口，使得不同研究者能够在较为一致的数据框架下比较算法效果。

围绕TCGA的分子分类任务，研究者最初主要采用传统统计学与机器学习方法，例如差异分析、单因素筛选、LASSO、SVM、随机森林等。这一阶段的核心目标是把高维组学数据压缩到可解释、可训练的子集，再在此基础上完成分类或生存分析。由于数据样本量通常不大、特征维度极高，这类方法长期以来仍然是最稳定的基线选择。

多组学融合研究一般可以划分为三条主线。第一条主线是早期融合，即把不同组学经过对齐和预处理后直接拼接到同一输入空间中。Concat属于这类方法的最简实现，优点是简单直接、易于比较，缺点是无法显式建模组学之间的交互。第二条主线是中期融合或表示学习，即先对每个组学提取低维表示，再在表示空间进行融合。MOFA、AE、VAE等方法都属于这一类，与直接拼接相比，这类方法能够显式压缩噪声、增强共享结构，并为后续解释提供潜在因子或嵌入向量。第三条主线是后期融合或决策级融合，即每种组学各自训练一个模型，再在输出层进行投票、加权或堆叠。后期融合更容易处理模态异构性，但往往难以利用跨模态共享结构，因此在强监督任务中不一定优于中期融合。

MOFA~\cite{mofa2018}是多组学研究中最具代表性的潜变量模型之一，核心思想是通过矩阵分解提取跨组学共享的低维潜在因子，同时保留每个组学的特有变化。与深度学习模型相比，MOFA的优势不在于大规模非线性表达能力，而在于可解释性与稳定性，非常适合样本较少的生物医学任务。在乳腺癌、卵巢癌、神经胶质瘤等TCGA数据上，MOFA常被用来做亚型分离、潜在结构发现和特征解释。很多研究会把MOFA先作为无监督表征学习步骤，再把输出的因子喂给分类器。本文采用的策略与这一思路一致，先学习低维因子，再用SVM完成分类，从而兼顾表示学习与分类性能。

近年来，多组学分类研究也逐渐引入图神经网络、对比学习、自监督预训练与生成式模型等方法。图神经网络的优势在于能显式建模基因-基因、样本-样本或通路-通路之间的关系，因此在某些任务中能够捕获单纯矩阵模型难以表达的结构信息。对比学习与自监督方法则更强调在标注有限时学习稳健表征，这对于临床数据稀缺场景很有吸引力。不过，这些新方法大多带来更高的实现复杂度，也更依赖图构建方式、负样本定义、预训练策略和大规模数据支持。相较之下，本文的目标不是追求模型谱系的最前沿，而是以一个可复现、可解释、便于课程报告展示的框架，系统比较传统多组学路线的实际效果。
